\documentclass[a4paper]{arrowhead}
\usepackage[T1]{fontenc}
\usepackage{longtable,booktabs,xurl,enumitem}
\setlength{\emergencystretch}{3em}
\setlength{\parindent}{0pt}\setlength{\parskip}{5pt}
\clubpenalty=10000\widowpenalty=10000
\setlist{nosep,leftmargin=*}
\usepackage{xurl}
\begin{document}
\ArrowheadTitle{LocalCloudManager / localCloudManagement-SD}\ArrowheadType{Service Description}\ArrowheadTypeShort{Rust}
\ArrowheadVersion{5.3.0}\ArrowheadDate{2026-11-09}
\ArrowheadAuthor{Eclipse Arrowhead migration project}\ArrowheadStatus{REVIEW DRAFT}
\ArrowheadContact{Project maintainers}\ArrowheadFooter{Eclipse Arrowhead -- Rust LocalCloudManager -- EPL-2.0}
\ArrowheadSetup
\section{Documentation edition and implementation scope}
This is the 5.3.0 documentation edition, dated 2026-11-09. It carries forward the system and service contracts from the retained 5.2.2 documentation and identifies the current primary implementation sources. The former \path|documentation/v5.2.2| folder is retained unchanged.

Document version 5.3.0 is distinct from executable, dependency, registered system/service and wire-contract versions. Values still defined by the source (including 5.2.2 artefact versions, 1.0.0 service versions and \path|SysMonitor-AMQP-IDD-5.2.2| payload identifiers) are preserved. This documentation update does not change runtime behaviour or establish a new implementation release. Historical verification statements describe the tests recorded in the inherited baseline; this revision validates document compilation and references.

\section{localCloudManagement service}
\par
Definition \path|localCloudManagement|, version1.0.0, producer LocalCloudManager, base \path|/api|. The actual source registers \path|generic_http|, protocol http, policy NONE. The console's password/session boundary is applied by the handlers, not expressed as certificate authentication by that descriptor.
\par
\subsection{Purpose and operations}
\par
The service exposes session login/state/logout, registry/runtime overview and grouping, historian data, consumers, intercloud summary, SysMonitor views and subscriptions, plant descriptions, orchestration rules, container logs and lifecycle. Read and mutation HTTP bindings are defined in the HTTP IDD. The registration descriptor contains only a subset of these routes; it must not be treated as the complete handler inventory.
\par
The operator authenticates with a scrypt-protected password and receives a random HttpOnly/SameSite=Strict cookie. Mutations require the explicit confirmation header. Missing password disables console APIs. A bearer identity token is not a substitute for an operator session.
\par
\subsection{Delimitations}
\par
Subscription mutation is acknowledgement-only, intercloud is a placeholder, and several displayed flags are synthetic. Provider calls use legacy declared SYSTEM identities and no native mTLS. This service is not certified interchangeable with the Java-cloud or newer Go/Rust support implementations. The relocation retains these semantics and documents them instead of claiming a security or parity upgrade.
\section{References and qualification}
See \path|docs/source-manifest.json| for original/relocated source SHA256 records, \path|VERIFICATION.md| for checks actually run, and \path|docs/LocalCloudManager.sysml| for the descriptive, unparsed variability model. No Java interchangeability, native TLS or sysMonitor publisher is claimed. Existing cloud source and running deployment remain unchanged.
\section{Revision record}
2026-11-09: Documentation edition 5.3.0 derived from the retained 5.2.2 sources; metadata and local references updated. Runtime and wire-contract versions remain source-defined.
\end{document}
